\documentclass[12pt]{article}

\usepackage{sbc-template}
\usepackage[brazil]{babel}
\usepackage{graphicx}
\usepackage{times}
\usepackage{url}

\sloppy

\title{Breve História da Psicologia Moderna Ocidental: Uma Leitura Pessoal sobre a Ciência, a Guerra e a Memória da Alma}

\author{Yago Melo da Costa}

\begin{document}
\maketitle

\begin{resumo}
A \emph{Breve História da Psicologia Moderna Ocidental}, de Teresa Sousa Machado (documento PDF, 128 p., lido integralmente; paginação impressa igual à do arquivo), é seguida aqui nome a nome, do laboratório de Leipzig às fronteiras de imigração dos Estados Unidos. Cada escola aparece como resposta a perguntas que o próprio século XX fabricou, e cada nome é ancorado em seu contexto de guerra, eugenia, exílio ou memória. A obra percorre as raízes filosóficas e fisiológicas da disciplina, a ascensão dos testes ao serviço da exclusão, o giro do vínculo e da memória coletiva, e o resgate de mulheres e disidentes que a história canônica silenciou.
\end{resumo}

\begin{abstract}
Teresa Sousa Machado's \emph{Brief History of Modern Western Psychology} (PDF document, 128 pp., read in full; printed pagination identical to the file's) is followed here name by name, from the Leipzig laboratory to the immigration frontiers of the United States. Each school appears as an answer to questions the twentieth century itself manufactured, and each name is anchored in its context of war, eugenics, exile, or memory. The work traces the discipline's philosophical and physiological roots, the rise of tests in the service of exclusion, the turn toward attachment and collective memory, and the recovery of women and dissenters whom canonical history silenced.
\end{abstract}

\section{O livro lido inteiro}

\noindent
\textbf{Edição utilizada:} Todas as citações referem-se ao documento PDF ``Breve História da Psicologia Moderna Ocidental'', com 128 páginas, assinado por Teresa Sousa Machado na primeira página \cite{machado}, lido integralmente para este ensaio. A paginação do arquivo corresponde à paginação impressa: o número da página está impresso no próprio texto, e todas as referências o seguem.

Uma epígrafe abre o livro, atribuída a Jacques Le Goff: a memória ``procura salvar o passado, para servir o presente e o futuro. Devemos trabalhar de forma a que a memória sirva para a libertação e não para a servidão dos homens'' \cite[p. 1]{machado}. A epígrafe funciona como programa: a \emph{Breve História} conta a psicologia como uma sequência de escolhas em que a memória oscila entre os dois destinos de Le Goff (libertação ou servidão), e o leitor é convocado a julgar teorias e condutas. A pergunta que organiza este artigo é, portanto, esta: o que o livro \emph{faz} com os nomes? E a resposta, que as seções seguintes demonstram, é a de que o livro faz com cada nome algo específico: recusa o vazio do ``grande homem'', ancora a teoria em seu contexto material e mostra, nome a nome, o que a disciplina escolheu e quem pagou por isso.

O material de partida já anuncia esse modo de contar. Contra a tentação de começos absolutos, o livro cita Correia Jesuíno: ``é sempre um pouco arbitrário assinalar começos absolutos em ciência, e a que psicologia não escapa à regra'' \cite[p. 1]{machado}. E desloca o eixo da narrativa: ``outras forças contextuais impulsionaram o desenvolvimento da nossa ciência, como as oportunidades económicas as guerras e a discriminação'' \cite[p. 1]{machado}. É uma história de ideias fabricadas \emph{por} séculos, não flutuando sobre eles. A frase de Hermann Ebbinghaus, repetida logo na primeira página, mede a escala do objeto: a psicologia tem ``um longo passado e uma curta história'' \cite[p. 1]{machado}. Passado longo, porque filosofia e fisiologia a antecipam por milênios; história curta, porque a autonomia cabe num século. E nessa curta história, o livro insiste, não há uma psicologia: ``a história da psicologia não existe, ao menos no singular'' \cite[p. 1]{machado}, com Cláudio Figueiredo, plural de partidos, de países, de feridas.

\section{Contextos antes de teorias}

O primeiro efeito da \emph{Breve História} é a recusa da neutralidade didática. O livro não alinha nomes em linha do tempo como quem arquiva; ele os lê como quem indaga. Wundt é mais que o homem do laboratório de Leipzig: é o ponto onde a fisiologia alemã decide separar-se da filosofia. Galton é mais que precursor dos testes: é o primo de Darwin que transforma parentesco em estatística. A escolha metodológica (contexto antes de teoria) tem consequência moral: quando a eugenia aparece, ela vem como conclusão de premissas que o livro expõe uma a uma, nunca como aberração inexplicável. O leitor vê a engrenagem funcionando antes de ver o estrago.

O registro linguístico do livro também faz parte do argumento. Escrito em português europeu (``raparigas'', ``investigação'', ``órfãos''), a \emph{Breve História} é uma leitura do cânone ocidental feita da margem atlântica da Europa: as escolas chegam de Leipzig, Viena, Londres, Cambridge, mas o olhar que as ordena é o de quem escreve perto de Lisboa e de Coimbra. Daí o duplo horizonte do livro: percorre o cânone universal (Wundt, Freud, Skinner, Piaget) e, ao mesmo tempo, devolve atenção às psicologias periféricas, a de Alves dos Santos e Sílvio Lima em Portugal, a de Helena Antipoff no Brasil. Uma história da psicologia escrita de dentro do centro contaria o centro; esta, escrita de fora, conta o centro e seus arredores, e os arredores, no fim, são onde se veem melhor as exclusões do centro. O resultado é um livro denso de nomes e sem herói único: a pluralidade é tese, não falha de síntese.

\section{As raízes: Wundt, Titchener e a primeira batalha das escolas}

O livro apresenta Wilhelm Wundt (1832--1920) como o ponto de autonomia da disciplina: ``o homem que estabeleceu definitivamente a autonomia da psicologia experimental foi incontestavelmente Wundt'' \cite[p. 10]{machado}, com o laboratório de Leipzig como marco. Wundt define o objeto com fórmula que o livro preserva: a psicologia é ``uma ciência empírica cujo objeto de estudo é a experiência interna, ou imediata'' \cite[p. 14]{machado}. E o mesmo Wundt é apresentado como precursor da \emph{Volkerpsychology}, ``a psicologia dos povos'', o estudo das produções culturais coletivas, linguagem, mito, costume \cite[p. 10]{machado}. O fundador, portanto, já carregava em si a tensão que definiria o século: laboratório de um lado, cultura do outro.

Edward Titchener (1867--1927), seu discípulo inglês, desenvolve o estruturalismo: analisar ``a consciência nos seus elementos mais básicos'', reproduzindo ``o método dos físicos'' \cite[p. 6]{machado}, com sensações, imagens e estados afetivos como elementos, introspecção treinada como instrumento, e a defesa de uma ciência pura, não aplicada \cite[pp. 22--23]{machado}. Em oposição, William James (1842--1910) e os funcionalistas americanos interessam-se pela função da mente no mundo real: para James, ``a consciência consistiria num fluxo contínuo de pensamentos'' \cite[p. 7]{machado}, irredutível à química elementar dos estruturalistas. O livro nomeia o desfecho com orgulho de família: o estruturalismo e o funcionalismo desencadearam as ``primeiras grandes batalhas das escolas de psicologia'' \cite[p. 6]{machado}.

Nasce assim o modelo de existência da disciplina: a escola como partido, a teoria como trincheira, a fundação como guerra de família. E o mesmo James reaparece mais tarde como o psicólogo que quase virou pacifista: em ``The Moral Equivalent of War'' (1906), reconhece os ``atrativos da guerra'' (lealdade, disciplina, honra) e propõe equivalentes morais para canalizá-los sem violência \cite[p. 100]{machado}. O livro registra o gesto porque ele pertence ao mesmo século que, trinta anos depois, poria a psicologia a serviço da guerra de verdade.

\section{Os testes, a eugenia e as guerras}

A parte mais sombria do livro, e também a mais necessária, é a dos testes psicológicos. Binet (1857--1911) e Simon (1873--1961) apresentam suas ideias no Congresso Internacional de Psicologia de Roma, em 1905, e em 1908 publicam ``uma escala métrica que avaliaria a inteligência através do teste Binet-Simon'', com propósito escolar explícito: saber ``em que classe se deveriam colocar'' as crianças ``para que tivessem o melhor aproveitamento na escolaridade'' \cite[p. 55]{machado}. Binet discordava de Galton e de Cattell: onde eles mediam processos sensório-motores, ele propunha tarefas de comparar, relacionar, afirmar e negar, e a inteligência valia como ``ciência da ação'' \cite[p. 55]{machado}.

Do outro lado do mesmo império das medidas está Francis Galton (1822--1911), primo de Darwin, que em ``O Gênio Hereditário'' (1869) compilou estudos de 300 famílias e legou a primeira análise quantitativa da inteligência \cite[p. 12]{machado}, e também a palavra ``eugenia'', com a pretensão de orientar o futuro do parentesco humano \cite[p. 26]{machado}. Nos Estados Unidos, Henry Goddard aplicou testes a imigrantes e produziu as tabelas que o livro não esconde: 83\% de baixo nível entre húngaros, 79\% entre italianos, e números semelhantes para russos, pesquisa que alimentou diretamente as leis de restrição à imigração dos anos 1920 \cite[p. 101]{machado}. O mesmo instrumento que Binet pensara para a sala de aula governou a fronteira. A régua que mede crianças mede povos: o teste, nesse trecho da história, é tecnologia de fronteira antes de ser tecnologia de cuidado.

A Primeira Guerra colocou a psicologia ``no mapa'' \cite[p. 100]{machado}; a Segunda a espalhou pela seleção, pela espionagem, pela propaganda, e o livro narra o resultado sem poupá-lo: em 1945, ``seis milhões de judeus tinham sido assassinados, e entre eles estariam mais de um milhão de crianças'' \cite[p. 95]{machado}, com relatos escritos ``à revelia'' que viriam a resgatar o que os campos tentaram apagar. As duas contas, 83\% e seis milhões, são os dois modos de contabilidade do mesmo século: a ciência que tabula inferioridades e a história que tabula cadáveres. O livro evita a acusação fácil de que a psicologia causou o Holocausto e prefere registrar algo mais difícil de administrar: que ela circulava confortavelmente entre as mesmas premissas de hierarquia que o tornaram possível.

\section{Psicanálise, vínculo e trauma: o giro da segunda metade}

A segunda metade do livro opera uma inversão de objeto: onde a primeira media, a segunda vincula. Sigmund Freud (1856--1939) aparece nas páginas do livro como o homem que usou ``pela primeira vez o termo psicanálise'' em 1896 e que, ao afirmar a influência do inconsciente e a ligação dos impulsos sexuais às neuroses, ``criou oposição nos meios científicos'' \cite[p. 34]{machado}. O livro registra o nascimento em Freiburg e a disputa que se seguiu, e preserva a imagem que atravessou o século: a mente como iceberg, consciente pequeno e visível sobre inconsciente vasto e propulsor \cite[p. 53]{machado}. Do mesmo campo Freudiano saem as duas herdeiras que o livro trata com atenção especial. Anna Freud (1895--1982) sistematiza os mecanismos de defesa do ego (repressão, negação, racionalização, projeção, entre outros) \cite[p. 39]{machado}, e Melanie Klein (1882--1960) inaugura a análise do brincar, exposta segundo o livro no título ``A Psicanálise das Crianças'' \cite[p. 53]{machado}, com o ``seio materno'' como primeiro objeto de uma relação marcada por estímulos ``oral-destrutivos'' \cite[p. 46]{machado}.

O giro do vínculo tem nome e páginas próprias. John Bowlby (1907--1990) é apresentado com a sentença que o livro transcreve em inglês antes de traduzir o efeito: a privação prolongada do cuidado materno ``may have grave and far-reaching effects on his character and so on the whole his future life'' \cite[p. 47]{machado}. E a fórmula que resume a ruptura com o behaviorismo da época: ``o amor da mãe pelo seu filho é tão vital para o desenvolvimento como a boa nutrição'' \cite[p. 47]{machado}. As fases da separação (protesto, desespero, desapego) \cite[p. 48]{machado} completam o quadro. O livro dedica a mesma atenção aos experimentos que forçaram o campo a rever seus dogmas: René Spitz (1887--1974) e a síndrome de hospitalismo das instituições \cite[p. 43]{machado}; Harry Harlow (1905--1981) e os macacos Rhesus que preferiam a ``mãe'' de pano felpudo ao arame que alimentava \cite[p. 52]{machado}; Donald Winnicott (1896--1971), que em 1947 destaca ``o efeito devastador das múltiplas mudanças de alojamento/acolhimento'' das crianças da Guerra, cujas privações ``tendem a tornar-se delinquentes'' \cite[p. 50]{machado}; e Mary Ainsworth (1913--1999), com a ``Strange Situation'' e os modelos internos de funcionamento que a criança levará por todo o ciclo vital \cite[pp. 109--110]{machado}.

O capítulo do vínculo desemboca na memória, e o livro escolhe seu teórico: Maurice Halbwachs (1877--1945) e a memória coletiva, fator de identidade dos grupos diante do trauma \cite[p. 97]{machado}. Diário escondido, fotografia salva, relato de sobrevivente, memória que resiste ao negacionismo. O livro faz consigo mesmo o que faz com a psicologia: salva o passado ``para servir o presente e o futuro'', conforme a epígrafe de Le Goff que abre o volume, e a história da disciplina vira memória coletiva da própria disciplina.

\section{As escolas como respostas ao século}

O livro apresenta estas escolas sem rivalidades de gabinete: cada uma atende a uma demanda do próprio século, e a leitura fica mais fácil quando se pergunta o que cada uma foi chamada para fazer. O behaviorismo de John B. Watson (1878--1958) é apresentado como a decisão de ``abandonar a preocupação com a mente'' e estudar apenas o comportamento \cite[p. 3]{machado}, e a fórmula que o livro preserva define o programa: o behaviorismo é ``um ramo puramente experimental da ciência natural'', cujo objetivo teórico é a previsão e o controle \cite[p. 82]{machado}. Por trás dessa escolha há um desejo de século: uma disciplina nova querendo a precisão da física e trocando o sujeito pelo experimento para consegui-la. Os alunos do método completaram o programa: Edward Thorndike (1874--1949) legou a lei do efeito, a resposta que se fortalece com a satisfação e se enfraquece com o desprazer \cite[p. 111]{machado}, e Ivan Pavlov (1849--1936) o condicionamento clássico do sino e da comida \cite[p. 113]{machado}, com a biografia que o livro conta com sabor: filho de padre que abandonou a carreira religiosa, e que em 1924, quando expulsaram os padres da Academia Médica Imperial, respondeu com renúncia, ``eu também sou filho de um padre'' \cite[p. 67]{machado}.

A Gestalt entra com a lição oposta, o todo anterior às partes. Wertheimer (1880--1943) formulou ``as leis básicas da perceção'' em 1923 \cite[p. 93]{machado}, e o livro guarda a cena do movimento aparente observado em 1910, num comboio da Renânia, duas lâmpadas alternando a ilusão de movimento contínuo \cite[p. 93]{machado}. A cena importa porque ela acontece dentro de um trem, e o trem, naquele século, organizava horários, guerras e migrações: a percepção humana descobre sua ilusão no veículo com que o século se movia.

O humanismo de Carl Rogers (1902--1987), a ``terceira força'', a terapia ``centrada no cliente'', onde o terapeuta facilita e o sujeito decide \cite[p. 81]{machado}, responde ao método que havia perdido o sujeito, e a hierarquia de Abraham Maslow (1908--1970), das necessidades fisiológicas à autorrealização \cite[p. 88]{machado}, devolve à psicologia uma seta para cima, para o que a pessoa pode se tornar. Jean Piaget (1896--1980) aparece com a psicologia genética da inteligência e os estádios do desenvolvimento, sustentados na recusa do ``pequeno adulto'': a criança constrói, não recebe \cite[pp. 74--75]{machado}.

E, no meio do percurso, o livro faz uma inserção que é, em si, um argumento de geografia: Paulo Freire (1921--1997). Os oprimidos que culpam ``causas mágicas e míticas'' (o destino, a sina, a vontade de Deus) e a autonomia como conquista, não como dom \cite[p. 77]{machado}; os 300 adultos alfabetizados em menos de 40 horas em 1963 \cite[p. 78]{machado}; a prisão e o exílio sob a ditadura. Freire é o único nome do capítulo das escolas que vem de fora do eixo europeu-americano, e o livro o coloca lá sem cerimônia, como quem abre a janela da síntese para o ar de fora.

\section{Mulheres, exílios e a psicologia em Portugal}

Uma das decisões mais legíveis do livro é o resgate nominal das mulheres que a história canônica tratou como nota de rodapé. Leta Hollingworth (1886--1939), cujas histórias não foram publicadas por ser mulher, chefe do laboratório psicológico do Hospital Central de Bellevue, nos anos 1920 ``passou a trabalhar com crianças sobredotadas'', desenvolvendo as estratégias que Terman apontara sem detalhar \cite[p. 70]{machado}. Dorothea Dix (1802--1887), que em 1841 percorreu dezoito meses de prisões, hospitais e asilos de Massachusetts e encontrou ``os doentes... tratados pior que os animais'', acorrentados e esquecidos, denúncia que reformou legislação \cite[p. 35]{machado}. Frances Power Cobbe (1822--1904), jornalista e ativista contra o espancamento das mulheres e contra a vivissecção, fundadora de sociedades de proteção animal e defensora do sufrágio \cite[pp. 26--27]{machado}. Helena Antipoff (1892--1974), russa de formação parisiense, que em 1929 veio ao convite de Minas Gerais para a ``Reforma Francisco Campos-Mário Cassanta'' \cite[p. 73]{machado}, fundou em Belo Horizonte o primeiro laboratório de psicologia aplicada da América do Sul e conduziu a Fazenda do Rosário, dedicada às crianças ``excecionais'' \cite[p. 89]{machado}.

A mesma atenção o livro devota à psicologia em Portugal, o que confirma sua localização: Alves dos Santos e o laboratório que fundou; Egas Moniz; Faria de Vasconcelos; e Sílvio Lima (1904--1993), o conimbricense introdutor da moderna psicologia portuguesa que ``em Maio de 1935 foi afastado do ensino por ser oposicionista ao Estado Novo'' \cite[p. 57]{machado}. A ciência sob o mesmo regime que a grande história europeia narraria noutro capítulo: a psicologia portuguesa consolida-se como ensino universitário público apenas ``após a revolução de 1974'' \cite[p. 57]{machado}. O leitor percebe então o desenho do livro inteiro: mulheres e disidentes são o contraponto que a síntese honesta obrigava a incluir, porque uma história da psicologia que ignore quem excluiu reproduz, no ato de contar, a exclusão que narrou.

\section{Os nomes que a régua excluiu}

Ler as 128 páginas de Teresa Sousa Machado foi atravessar o século XX pelas oficinas onde as teorias foram forjadas: laboratórios, campos de guerra, orfanatos, fronteiras de imigração. Foi, para os meus propósitos de 2026, o livro que faltava entre a clínica que quero exercer e as literaturas que venho lendo: ele dá o passado que o \emph{Procognitiva} \cite{procognitiva2014} pressupõe sem contar.

Com Kafka \cite{kafka2022}, a conexão é a contabilidade. As tabelas de Goddard (83\%, 79\%) são o registro da mesma lógica que reduz Gregor Samsa a um problema a ser processado pela burocracia doméstica: mede-se, tabula-se, e o que não cabe na tabela deixa de ser pessoa. A diferença de escala é a única diferença: a mesa de exame de imigrantes é o quarto de Gregor em tamanho de nação. Com Krenak \cite{krenak2019}, a conexão é o conceito de humanidade: a ``humanidade esclarecida'' que, segundo o livro de Krenak, saía ``colonizando o resto do mundo'' \cite[p. 11]{krenak2019} encontrou, nos testes de Goddard, a tabela que faltava para separar quem era gente de quem não era, e o ``liquidificador chamado humanidade'' \cite[p. 14]{krenak2019} tem, no Binet-Simon desviado de seu propósito escolar, um dos primeiros motores. A eugenia é o lugar onde Krenak e Machado se reconhecem: o passo anterior à extração é sempre a desclassificação de alguém como gente.

Com Saramago \cite{saramago1997}, a conexão é o mapa: a frase de Figueiredo, ``a história da psicologia não existe, ao menos no singular'', é a versão disciplinar da descoberta saramaguiana de que ``estão todas nos mapas'' só enquanto houver quem desenhe os mapas. Cada história da psicologia é uma ilha pintada por alguém em algum lugar; a de Machado foi pintada perto de Coimbra, e é por isso que nela se vê, com clareza rara, Portugal e o Brasil recebendo a mesma atenção que Leipzig e Viena. E com a amostra de duas páginas do \emph{Procognitiva} \cite{procognitiva2014}, que li antes deste livro, o círculo se fecha: o programa de atualização é a resposta institucional ao problema que Machado levanta, para que ``a curta história'' não vire amnésia, publica-se um capítulo por mês; a pergunta que sobra, e que o muro de acesso da amostra materializa, é de quem é o direito à memória quando ela vira assinatura.

Fecho com as palavras que abrem o livro, porque são as que ficaram: a memória deve ``servir para a libertação e não para a servidão dos homens'' \cite[p. 1]{machado}. Esta \emph{Breve História} cumpre o que promete, salva o passado da psicologia para servir o presente de quem vai praticá-la, e me deixa com uma obrigação modesta e duradoura: contar a curta história sem encolher o longo passado, lembrar os nomes que ficaram de fora das tabelas, e tratar a alma como parente, nunca como medida.

\begin{thebibliography}{99}

\bibitem{machado}
MACHADO, Teresa Sousa.
\emph{Breve História da Psicologia Moderna Ocidental}.
Documento PDF, 128 p.

\bibitem{krenak2019}
KRENAK, Ailton.
\emph{Ideias para adiar o fim do mundo}.
1ª ed. São Paulo: Companhia das Letras, 2019.
ISBN 978-85-359-3241-6.

\bibitem{kafka2022}
KAFKA, Franz.
\emph{A Metamorfose} [\emph{Die Verwandlung}, 1915].
Tradução de Giovane Rodrigues.
São Paulo: Instituto Mojo, 2022. (Coleção Literatura Livre).
Edição bilíngue Português-Alemão. E-book.
ISBN 978-65-89008-23-1.
Disponível em: \url{https://mojo.org.br/wp-content/uploads/2024/08/006-A-metamorfose-978-65-89008-23-1.pdf}.

\bibitem{saramago1997}
SARAMAGO, José.
\emph{O Conto da Ilha Desconhecida}.
São Paulo: Companhia das Letras, 1997.

\bibitem{procognitiva2014}
NEUFELD, Carmem Beatriz; FALCONE, Eliane Mary de Oliveira; RANGÉ, Bernard (Orgs.).
\emph{Procognitiva: Programa de Atualização em Terapia Cognitivo-Comportamental. Ciclo 1, Volume 1}.
Porto Alegre: Artmed Panamericana, 2014. 151 p.

\end{thebibliography}

\end{document}